\chapter{認識 LaTeX：從排版問題開始}

本章先從數位排版的需求出發，認識 TeX 的起源、LaTeX 所扮演的角色，並釐清初學者常遇到的幾個疑問。讀完本章後，你不必背誦歷史細節；更重要的是理解「內容」與「排版規則」如何分工，以及什麼情況適合使用 LaTeX。本章的概念參考了李果正的講義\cite{lee2004}。

\section{一套為精確排版而生的系統}

\subsection{Knuth、排版與 TeX}

TeX 的誕生和一本大型計算機科學著作的排版經驗有關。美國計算機科學家 Donald E. Knuth 在整理《The Art of Computer Programming》（常縮寫為 TAOCP）時，對當時印刷成品的數學公式與版面品質不滿意，於是著手設計一套能細緻控制字距、行距與符號位置的排版系統。TeX 的開發始於 1970 年代末，後續版本逐步成熟，成為學術出版與技術文件的重要工具。

TeX 的特色不只是「把文字變成 PDF」。它會依據文字、字型與版面規則計算斷行和頁面配置；數學公式、表格或特殊符號也能按照一致的規則排版。作者可以描述要呈現的內容與結構，再由排版系統處理大量細節。

TAOCP 是 Knuth 長期撰寫的演算法與程式設計著作。它在這段歷史中扮演的是催化劑：大型、含大量數學內容的書籍，讓精確排版問題變得格外明顯。學習 LaTeX 並不需要先讀完 TAOCP，也不需要先成為數學或計算機專家。

\subsection{從精密控制到方便寫作}

TeX 本身提供非常有彈性的低階控制能力；若每篇文章都直接從細節開始指定，寫作會變得繁瑣。LaTeX 建立在 TeX 之上，提供較高階的文件結構與命令，例如章節、標題、清單、交互參照與參考文獻。作者可以寫出「這是章節標題」或「這裡引用一筆文獻」，再由文件類別和套件決定整體呈現方式。

可以用分工來理解兩者：TeX 提供排版引擎和底層能力，LaTeX 則提供一套常用的文件規則，讓作者較容易組織內容。這並不表示排版細節永遠不能調整，而是先有合理的一致預設，需要時再逐步修改。

\section{LaTeX 文件如何變成成品}

LaTeX 文件通常以純文字的 \texttt{.tex} 檔保存。編輯器負責撰寫原始碼，XeLaTeX 等編譯程式讀取原始碼、套用排版規則，最後產生 PDF。編譯過程也可能建立輔助檔案，用來記錄目錄、頁碼或交互參照等資訊。

\begin{center}
原始碼（.tex）\quad $\longrightarrow$\quad XeLaTeX 編譯\quad $\longrightarrow$\quad PDF
\end{center}

這種工作方式和直接在頁面上拖曳文字方塊不同。原始碼中可以清楚看出標題層級、段落、公式及引用關係；同一份內容也能套用不同的文件類別或樣式。文件越長、結構越明確，這種分離通常越有幫助。

以下是一個精簡的繁體中文文件骨架。它使用 XeLaTeX，並以 xeCJK 設定中文排版；字型名稱必須是目前電腦已安裝、且 XeLaTeX 能找到的字型。

\begin{verbatim}
\documentclass[a4paper,12pt]{article}
\usepackage{fontspec}
\usepackage{xeCJK}
\setCJKmainfont{DFKai-SB}
\setmainfont{Times New Roman}
\XeTeXlinebreaklocale "zh"
\XeTeXlinebreakskip = 0pt plus 1pt

\title{第一份中文文件}
\author{作者姓名}
\date{\today}

\begin{document}
\maketitle
這是一段繁體中文。
\end{document}
\end{verbatim}

這個例子展示的是文件的基本組成，不是唯一寫法。不同課程、期刊或出版單位可能會使用不同的文件類別與套件；只要掌握內容結構，就能再依需求調整外觀。

\section{初學者常見的幾個疑問}

\subsection{一定要使用特定作業系統嗎？}

不需要。TeX 發行版與 LaTeX 編輯工具可在 Windows、macOS 與 Linux 等平台使用，也有線上編輯與編譯服務。不同平台的安裝步驟、字型名稱和可用套件可能不同，但文件原始碼通常可以在平台之間移動。若文件指定了特定字型，其他電腦也必須有該字型，或由作者改用可用的替代字型。

\subsection{LaTeX 已經過時了嗎？}

不能只用「新」或「舊」判斷工具是否合適。LaTeX 仍持續維護，並廣泛用於數學、科學、工程及技術文件，也能用於書籍、講義和一般長篇文件。另一方面，簡短且以視覺拖曳為主的文件，使用一般文書軟體可能更直接。選擇工具時，應考慮文件結構、公式需求、協作方式、出版規格與維護成本。

\subsection{LaTeX 是所見即所得嗎？}

傳統工作流程通常是先編輯原始碼，再編譯並檢視 PDF，因此不是每個字元一輸入就立刻呈現在最終頁面上。這種方式常被概括為「所想即所得」：作者描述文件的結構和內容，排版系統依規則產生結果。現代編輯器可以同時顯示預覽，讓修改與檢查更方便；不過預覽並不改變「編輯來源、編譯成品」的基本流程。

\subsection{它很難學，而且只能排數學嗎？}

LaTeX 有自己的命令語法，因此剛開始需要熟悉反斜線、括號、環境與編譯錯誤訊息。初期不必追求一次學會所有套件；從標題、段落、清單、公式和引用等常見需求開始，逐步累積即可。它以數學排版聞名，但並不限於數學文章，也可用於報告、講義、書籍與結構化技術文件。

\section{本課程從哪裡開始}

本課程會先練習建立文件、編譯與檢查 PDF，接著逐步加入標題、段落、清單、表格、圖片與公式。遇到中文時，除了原始碼內容，也要注意編譯引擎、字型是否安裝，以及中文斷行設定。這些因素會共同影響最終輸出。

初學階段可以先建立三個習慣：
\begin{enumerate}
    \item 將內容與格式命令分清楚，讓文件結構容易閱讀。
    \item 每次完成一小段修改就編譯，依錯誤訊息定位問題。
    \item 檢查 PDF 成品，而不只確認原始碼看起來合理。
\end{enumerate}

\subsection*{本章回顧}

TeX 回應的是精確排版的需求；LaTeX 則在 TeX 的能力之上提供較方便的文件結構。作者以原始碼描述內容，編譯程式依規則產生 PDF。LaTeX 是否適合某份文件，取決於文件的結構與工作流程，而不是單看它是否以公式排版聞名。

\subsection*{思考與討論}

\begin{enumerate}
    \item TeX 與 LaTeX 分別負責什麼？請用自己的話說明兩者的分工。
    \item 如果一份文件需要大量公式、章節與交互參照，使用 LaTeX 可能有哪些好處？
    \item 在另一台電腦編譯含中文的文件時，除了原始碼之外，還應確認哪些環境條件？
\end{enumerate}

\clearpage
\addcontentsline{toc}{chapter}{\bibname}
\begin{thebibliography}{9}
\bibitem{lee2004}
李果正（Edward G. J. Lee），\textit{大家來學 \LaTeX{}}，Version 1.0，2004 年 3 月 8 日。本文僅以該講義作為概念參考，內容與例子另行撰寫。
\end{thebibliography}